% reality/coordinate_systems.tex

\chapter{Coordinate Systems and World Embedding}

\label{chap:coordinate-systems}

This chapter owns the coordinate and CRS explanation of the Reality
part. In AIM, CRS and coordinate systems define embedding and
transformation context; they do not define constructive alignment
identity.

\section*{CRS is not Metadata}

\begin{insight}
In many software systems, the coordinate reference system (CRS) is
treated as metadata.

It is stored, displayed, and occasionally converted — but rarely
considered part of the model itself.

\medskip

In railway alignment modelling, this view is fundamentally wrong.

\medskip

The CRS determines:
\begin{itemize}
\item how distances are measured,
\item how curvature is interpreted,
\item and how different datasets relate to each other.
\end{itemize}

\medskip

A shift of a few centimeters in coordinates may appear negligible.

A change in curvature is not.

\medskip

If coordinate systems are inconsistent:
\begin{itemize}
\item alignments no longer match,
\item curvature estimates become unstable,
\item and optimization results lose physical meaning.
\end{itemize}

\medskip

The problem is not numerical.

It is structural.

\medskip

A railway alignment does not exist independently of its spatial
reference.

The coordinate system is part of the model.
\end{insight}


% ============================================================
\section{Coordinate Reference Systems}
% ============================================================

\begin{definition}[Coordinate reference system]
A coordinate reference system (CRS) is a mathematical framework that
maps spatial coordinates to positions in a defined world embedding.
\end{definition}

A CRS combines reference surface assumptions, coordinate definitions,
and transformation conventions.

% ============================================================
\section{Embedding Context}
% ============================================================

Railway projects commonly use multiple systems simultaneously, including
ellipsoid-referenced systems, projected systems, and local engineering
frames. These systems provide embedding context for geometry and
measurements.

\begin{principle}[CRS context principle]
CRS belongs to realization and embedding context, not to constructive
identity.
\end{principle}

% ============================================================
\section{Transformations and Consistency}
% ============================================================

\begin{definition}[Coordinate transformation]
A coordinate transformation is a mapping
\[
T:\R^n\rightarrow\R^n
\]
between coordinate systems.
\end{definition}

Transformation handling must be explicit and auditable, since small
transformation inconsistencies can create large interpretation errors in
curvature-related or projection-related quantities.

% ============================================================
\section{Local Frames and Numerical Stability}
% ============================================================

\begin{definition}[Local frame]
A local frame is a coordinate system anchored near an engineering region
of interest.
\end{definition}

Local frames improve numerical stability for runtime evaluation,
projection, and reconstruction, while preserving linkage to global
reference context.

% ============================================================
\section{Stationing versus CRS Coordinates}
% ============================================================

CRS coordinates express embedded position. Stationing expresses route
address along railway reference structures. These are distinct layers
and must not be conflated.

% ============================================================
\section{Relation to Metric Realization}
% ============================================================

Coordinate handling consumes metric realization
(\cref{chap:metric-realization}) and realization operators; it does not
replace them. World embedding and world--track relations remain
realization-dependent services.

% ============================================================
\section{Connection to AIM}
% ============================================================

\begin{summary}
Within AIM, coordinate systems and CRS are explicit embedding context
for realized and observed railway states.

They are essential for consistency and interoperability, but they are
not alignment identity.
\end{summary}
